\documentclass{article}
%\setlength{\textwidth}{15cm}

%\documentclass[submit]{aiaa-tc}% insert '[draft]' option to show overfull boxes

\usepackage[margin=1.3in]{geometry}
 \usepackage{wrapfig}% embedding figures/tables in text (i.e., Galileo style)
 \usepackage{threeparttable}% tables with footnotes
 \usepackage{dcolumn}% decimal-aligned tabular math columns
  \newcolumntype{d}{D{.}{.}{-1}}
 \usepackage{nomencl}% automatic nomenclature generation via makeindex
  \makeglossary
 \usepackage{subfigure}% subcaptions for subfigures
 \usepackage{subfigmat}% matrices of similar subfigures, aka small mulitples
 \usepackage{fancyvrb}% extended verbatim environments
  \fvset{fontsize=\footnotesize,xleftmargin=2em}
 \usepackage{lettrine}% dropped capital at beginning of paragraph
 \usepackage[dvips]{dropping}% alternative dropped capital package
 \usepackage{hyperref}% embedding hyperlinks [must be loaded after dropping]
 \usepackage{amsmath}
 \usepackage{color}
 \usepackage{setspace}
 \usepackage{epsf}
 \usepackage{epsfig}
% \renewcommand{\baselinestretch}{1.25}

\begin{document}

\begin{titlepage}
\begin{center}

% Upper part of the page. The '~' is needed because \\
% only works if a paragraph has started.
\includegraphics[width=0.25\textwidth]{figures/LLlogo.eps}~\\[1cm]

\textsc{\LARGE LiberLocus}\\[1.5cm]

% Title
{ \huge \bfseries LLVisu Library}\\[0.4cm]

\textsc{\Large LiberLocus scientific output generation library for serial software}\\[0.5cm]

\textsc{Version 1.4}\\[3.5cm]

\textsc{\Large LiberLocus Team}\\[1.5cm]

\textsc{\Large Earth}\\[0.5cm]

\vfill

% Bottom of the page
{\large \today}

\end{center}
\end{titlepage}

%\maketitle
%\section{Introduction}

\section{Introduction}
\label{sec:intro}

This is a documentation about the LiberLocus LLVisu library. LLVisu aims to bring all the scientific data formats under a single library. LLVisu is written in OOP (in C++) in order to facilitate newly emerging data formats very easily through definition of new objects.

So far LLVisu supports only unstructured meshes. But considering the fact that structured meshes are a subset of unstructured meshes, one can get an output for a structured mesh using LLVisu by defining the regular cell connectivities. 

The supported file formats are VTK\cite{VTK} and XDMF\cite{XDMF} (through HDF5\cite{HDF5}). User can generate output data with one or multiple types simultaneously with the appropriate function calls.

The installation procedure is completed with CMake\cite{CMake}. Once installed, by including the appropriate header file in your code, functions can be called for output generation. Detailed explanation of these functions can be found in Section~\ref{sec:API}

%\section{What is new in Version 2.0}
%\label{sec:whatsnew}

\newpage

\section{LLVisu API}
\label{sec:API}

\noindent {\large void {\bf liberVisuVTK}(\newline \par {\it string}  fileName, {\it int}  nodeNum, {\it float}  *pos\_x, {\it float}  *pos\_y, {\it float}  *pos\_z, {\it int}  cellNum, {\it int}  nodePerCell, {\it int}  **cells, {\it char}  **varName, {\it int}  **varSize, {\it float}  **varMatrix \newline) }
\bigskip

\noindent {\large {\bf Description}}
        \par This function creates a vtk type output file in vtu format.
\bigskip
        
\noindent {\large {\bf Parameters}}
        \par Below are the definitions of the parameters in the function.
\medskip

        {\bf fileName}  \par the absolute path to the output directory + filename
\medskip
        
        {\bf nodeNum} \par number of nodes present in the mesh
\medskip
        
        {\bf pos\_x} \par pointer to the array that stores the x position of the nodes
\medskip
        
        {\bf pos\_y} \par pointer to the array that stores the y position of the nodes
\medskip
        
        {\bf pos\_z} \par pointer to the array that stores the z position of the nodes
\medskip
        
        {\bf cellNum} \par number of cells present in the mesh
\medskip
        
        {\bf nodePerCell} \par number of nodes per cells 
\medskip
        
        {\bf cells} \par pointer to the two dimensional array that stores the cell connectivity
\medskip
        
        {\bf varName} \par pointer to the array that stores the variable names
\medskip
        
        {\bf varSize} \par number of variables that are passed 
\medskip
        
        {\bf varMatrix} \par pointer to the two dimensional array that stores the variable value arrays(cell or node based) 

\newpage

\noindent {\large void {\bf liberVisuXMF}(\newline \par {\it string}  fileName, {\it string}  gridName, {\it string}  topoType, {\it int}  nodeNum, {\it float}  *pos\_x, {\it float}  *pos\_y, {\it float}  *pos\_z, {\it int}  cellNum, {\it int}  nodePerCell, {\it int}  **cells, {\it char}  **varName, {\it char}  **varType, {\it int}  *varLength, {\it int}  varSize, {\it float}  **varMatrix \newline) }
\bigskip

\noindent {\large {\bf Description}}
        \par This function creates a xdmf type output file in xmf format using hdf5.
\bigskip
        
\noindent {\large {\bf Parameters}}
        \par Below are the definitions of the parameters in the function.
\medskip

        {\bf fileName}  \par the absolute path to the output directory + filename
\medskip
        
        {\bf nodeNum} \par number of nodes present in the mesh
\medskip
        
        {\bf pos\_x} \par pointer to the array that stores the x position of the nodes
\medskip
        
        {\bf pos\_y} \par pointer to the array that stores the y position of the nodes
\medskip
        
        {\bf pos\_z} \par pointer to the array that stores the z position of the nodes
\medskip
        
        {\bf cellNum} \par number of cells present in the mesh
\medskip
        
        {\bf nodePerCell} \par number of nodes per cells 
\medskip
        
        {\bf cells} \par pointer to the two dimensional array that stores the cell connectivity
\medskip
        
        {\bf varName} \par pointer to the array that stores the variable names
\medskip
        
        {\bf varSize} \par number of variables that are passed 
\medskip
        
        {\bf varMatrix} \par pointer to the two dimensional array that stores the variable value arrays(cell or node based) 

        


  \begin{thebibliography}{1}

  \bibitem{VTK} Visualization Tool Kit, Page {\em http://www.vtk.org/}
  \bibitem{XDMF} Extensible Data Model Format, Page {\em http://www.xdmf.org}
  \bibitem{HDF5} Hierarchical Data Format Group, Page {\em http://www.hdfgroup.org/HDF5/}
  \bibitem{CMake} Cross-platform, open-source build system, Page {\em http://www.cmake.org/}
%  \bibitem{wiki2} Mergesort Wikipedia Page, {\em http://en.wikipedia.org/wiki/Merge\_sort}
%  \bibitem{main} Atanas Radenski "{\em Shared Memory, Message Passing, and Hybrid Merge Sorts for Standalone and Clustered SMPs}"  International Conference on Parallel and Distributed Processing Techniques and Applications, 2011 pp. 367-373.
%  \bibitem{sec} Jaeyoung Park, Kyong-Gun Lee and Jong Tae Kim "{\em Parallel Mergesort Implementation using OpenMP}" Proceedings of The 2011 World Congress in Computer Science, Computer Engineering and Applied Computing. "http://world-comp.org/p2011/PDP8400.pdf"
%
  \end{thebibliography}

\end{document}
